\documentclass[11pt,a4paper]{article}

\usepackage{amsmath, amssymb, amsthm}
\usepackage{geometry}
\usepackage{booktabs}
\usepackage{hyperref}
\usepackage{microtype}
\usepackage{enumitem}
\usepackage{xcolor}

\geometry{margin=2.5cm}

\theoremstyle{plain}
\newtheorem{theorem}{Theorem}[section]
\newtheorem{proposition}[theorem]{Proposition}
\newtheorem{corollary}[theorem]{Corollary}

\theoremstyle{definition}
\newtheorem{definition}[theorem]{Definition}
\newtheorem{example}[theorem]{Example}

\theoremstyle{remark}
\newtheorem{remark}[theorem]{Remark}

% Shorthand
\newcommand{\GH}{\mathrm{GH}(2,2)}
\newcommand{\Gtwo}{G_2(2)}
\newcommand{\PSU}{\mathrm{PSU}(3,3)}
\newcommand{\PSL}{\mathrm{PSL}}
\newcommand{\Aut}{\mathrm{Aut}}
\newcommand{\F}{\mathbb{F}}
\newcommand{\Z}{\mathbb{Z}}

\title{\textbf{Spectral Theory of the Exceptional Lie Group $G_2(2)$}\\
\large Hexagonal Geometries and Their Graphs}
\author{Robin Langer}
\date{\today}

\begin{document}
\maketitle

\begin{abstract}
The exceptional Lie group $G_2(2)$ is the automorphism group of the split Cayley hexagon
$\mathrm{GH}(2,2)$, a generalized hexagon of order $(2,2)$ arising from the split octonion
algebra over $\mathbb{F}_2$.  We study three graphs whose automorphism group is $G_2(2)$:
the Tutte 12-cage (the incidence graph of $\mathrm{GH}(2,2)$, 126 vertices, cubic),
the $\mathrm{GH}(2,2)$ collinearity graph (63 vertices, 6-regular, distance-regular),
and the Meinhold-72 graph (72 vertices, 7-regular, girth 4).
We compute the full adjacency spectrum of each graph and analyse their Ramanujan status,
algebraic connectivity, and bipartiteness.  Despite sharing the same automorphism group of
order 12096, the three graphs exhibit strikingly different spectral profiles: the collinearity
graph is Ramanujan with only four distinct eigenvalues, the Tutte cage is bipartite with a
symmetric spectrum containing irrational eigenvalues $\pm\sqrt{6}$ and $\pm\sqrt{2}$,
and the Meinhold-72 graph is a bipartite strongly regular graph in a suitable sense with
eigenvalues $\pm 7$, $\pm\sqrt{13}$, and $\pm 1$.
\end{abstract}

\tableofcontents

\section{Introduction}

Symmetry imposes rigid constraints on spectra.  A vertex-transitive graph on $n$ vertices
with automorphism group $G$ has its adjacency spectrum determined by the representation
theory of $G$, and the most highly symmetric graphs often exhibit the cleanest spectral
profiles.  The exceptional algebraic group $G_2$ over finite fields provides some of the
most beautiful examples.

The split Cayley hexagon $\GH$ is a point-line geometry of order $(2,2)$: 63 points,
63 lines, every point on 3 lines, every line through 3 points, and the incidence graph
has girth 12.  Its full automorphism group is $\Gtwo$, a finite group of order 12096
whose derived subgroup is $\PSU$ of order 6048.  From $\GH$ one naturally extracts
three graphs, all sharing the same automorphism group but with very different combinatorial
and spectral properties:

\begin{enumerate}
  \item The \emph{Tutte 12-cage}: the point-line incidence graph of $\GH$, on 126 vertices,
        3-regular, of girth 12.
  \item The \emph{$\GH$ collinearity graph}: the graph on the 63 points where two points
        are adjacent if and only if they lie on a common line, yielding a 6-regular
        distance-regular graph with intersection array $\{6,4,4;1,1,3\}$.
  \item The \emph{Meinhold-72 graph}: a quotient of the Clayworth graph by $K = \PSL(3,2)$
        acting on 72 vertices, 7-regular, of girth 4 and diameter 5.
\end{enumerate}

In this paper we compute the full adjacency spectrum of each graph and draw out the
structural lessons encoded in those spectra.

\section{The Group $G_2(2)$ and Split Octonions}

\subsection{Octonions over $\mathbb{F}_2$}

The split octonion algebra $\mathbb{O}'$ over $\mathbb{F}_2$ is an 8-dimensional algebra
over $\mathbb{F}_2$ with a non-degenerate quadratic form $N$ of Witt index 4.  The group
of $\mathbb{F}_2$-algebra automorphisms of $\mathbb{O}'$ is the exceptional algebraic group
$G_2$, and over the two-element field this yields the finite group $G_2(2)$ of order
\[
  |G_2(2)| = 2^6 \cdot 3^3 \cdot 7 = 12096.
\]
The derived subgroup $G_2(2)'$ is $\PSU(3,3)$ of order 6048, which is a simple group.
Thus $G_2(2)$ is an almost-simple group with simple socle $\PSU(3,3)$.

\subsection{The Split Cayley Hexagon $\GH$}

A \emph{generalized hexagon of order $(s,t)$} is a point-line geometry in which every
line contains $s+1$ points, every point is on $t+1$ lines, and the point-line incidence
graph has girth 12 (so there are no ordinary $k$-gons for $k < 6$ in the geometry).
The split Cayley hexagon $\GH$ has $s = t = 2$, giving:
\begin{itemize}[noitemsep]
  \item $1 + 2(1 + 2 + 4 + 8 + 16 + 32) = 63$ points,
  \item $63$ lines,
  \item each point on 3 lines, each line through 3 points.
\end{itemize}
The points lie on the parabolic quadric $Q(6,2) \subset \mathrm{PG}(6,2)$, and the lines
are certain totally isotropic lines with respect to the associated alternating trilinear
form coming from the octonion multiplication.  The full automorphism group of $\GH$
is $G_2(2)$, which acts regularly on ordered flags (incident point-line pairs).

\section{The Three Graphs}

\subsection{The Tutte 12-Cage}

\begin{definition}
The \emph{Tutte 12-cage} is the point-line incidence graph of $\GH$: its 126 vertices
are the 63 points and 63 lines of $\GH$, and a point $p$ is adjacent to a line $\ell$
if and only if $p \in \ell$.
\end{definition}

The Tutte 12-cage is the unique $(3,12)$-cage: the smallest 3-regular graph of girth 12.
It is distance-transitive with diameter 6, and $\Aut(\Gamma) \cong G_2(2)$.  As a
bipartite graph (points versus lines), its spectrum is symmetric about 0.

\subsection{The $\GH$ Collinearity Graph}

\begin{definition}
The \emph{$\GH$ collinearity graph} (or \emph{point graph}) is the graph on the 63 points
of $\GH$ in which two points are adjacent if and only if they are collinear (lie on a
common line of $\GH$).
\end{definition}

Since each point lies on 3 lines each containing 2 other points, the collinearity graph
is 6-regular.  It is strongly regular but not with the parameters of a srg; more precisely
it is distance-regular with diameter 3 and intersection array
\[
  \{6, 4, 4; 1, 1, 3\}.
\]
The distance distribution from any vertex $v$ is: 1 vertex at distance 0 ($v$ itself),
6 at distance 1, 24 at distance 2, and 32 at distance 3.

\subsection{The Meinhold-72 Graph}

\begin{definition}
The \emph{Meinhold-72 graph} is the quotient of the Clayworth graph by the action of
$K = \PSL(3,2) \cong \PSL(2,7)$, a subgroup of $G_2(2)$ of order 168.  It has 72 vertices,
is 7-regular, has girth 4, and diameter 5.
\end{definition}

The Clayworth graph is the Cayley graph $\mathrm{Cay}(G_2(2), S)$ for a suitable
generating set $S$; collapsing the orbits of $K$ yields Meinhold-72.  The graph was
named after Nick Meinhold who first identified it as a distinct object in the census
of cubic arc-transitive graphs related to $G_2(2)$.  Unlike the other two, Meinhold-72
is not distance-regular.

\section{Spectral Analysis}

Let $\Gamma$ be a $k$-regular graph on $n$ vertices with adjacency matrix $A$.
The eigenvalues of $A$ satisfy $k = \lambda_1 \geq \lambda_2 \geq \cdots \geq \lambda_n$.
A connected $k$-regular graph is \emph{Ramanujan} if
\[
  \max(\lambda_2, |\lambda_n|) \leq 2\sqrt{k-1}.
\]
Ramanujan graphs are optimal expanders: they achieve the spectral gap predicted by
the Alon--Boppana theorem.

\subsection{Tutte 12-Cage}

The Tutte 12-cage ($n = 126$, $k = 3$) has Ramanujan threshold $2\sqrt{2} \approx 2.828$.

\begin{table}[h]
\centering
\caption{Adjacency spectrum of the Tutte 12-cage ($n=126$, $k=3$, girth 12).}
\label{tab:tutte}
\begin{tabular}{@{}rrl@{}}
\toprule
Eigenvalue & Multiplicity & Exact form \\
\midrule
$3$        & $1$  & $k$ (trivial) \\
$\sqrt{6}$ & $21$ & $\approx 2.4495$ \\
$\sqrt{2}$ & $27$ & $\approx 1.4142$ \\
$0$        & $28$ & \\
$-\sqrt{2}$& $27$ & \\
$-\sqrt{6}$& $21$ & \\
$-3$       & $1$  & \\
\bottomrule
\end{tabular}
\end{table}

\begin{proposition}
The Tutte 12-cage has spectrum
\[
  \bigl\{3^1,\; (\sqrt{6})^{21},\; (\sqrt{2})^{27},\; 0^{28},\;
  (-\sqrt{2})^{27},\; (-\sqrt{6})^{21},\; (-3)^1\bigr\}.
\]
Its spectrum is symmetric about 0, consistent with bipartiteness.
It is \emph{not} Ramanujan: $\sqrt{6} \approx 2.449$ falls below the threshold
$2\sqrt{2} \approx 2.828$, so it narrowly misses. The spectral gap is $3 - \sqrt{6} \approx 0.551$.
\end{proposition}

The symmetry of the spectrum about 0 is a consequence of bipartiteness: for any bipartite
graph, if $\lambda$ is an eigenvalue then $-\lambda$ is also an eigenvalue with the same
multiplicity.  The Tutte cage is indeed bipartite (points vs.\ lines), so this is expected.

The 7 distinct eigenvalues match the diameter plus 1 ($\mathrm{diam} = 6$), confirming
that the Tutte cage is distance-regular (as expected from its distance-transitivity).
In a distance-regular graph, the number of distinct eigenvalues equals the diameter plus one.

\subsection{GH(2,2) Collinearity Graph}

The collinearity graph ($n = 63$, $k = 6$) has Ramanujan threshold $2\sqrt{5} \approx 4.472$.

\begin{table}[h]
\centering
\caption{Adjacency spectrum of the $\GH$ collinearity graph ($n=63$, $k=6$, girth 3).}
\label{tab:gh22}
\begin{tabular}{@{}rrl@{}}
\toprule
Eigenvalue & Multiplicity & Source \\
\midrule
$6$   & $1$  & trivial eigenvalue \\
$3$   & $21$ & \\
$-1$  & $27$ & \\
$-3$  & $14$ & \\
\bottomrule
\end{tabular}
\end{table}

\begin{proposition}
The $\GH$ collinearity graph has spectrum $\{6^1, 3^{21}, (-1)^{27}, (-3)^{14}\}$.
It has exactly 4 distinct eigenvalues, consistent with its diameter of 3.
It is \textbf{Ramanujan}: $\lambda_2 = 3 < 2\sqrt{5} \approx 4.472$.
The spectral gap is $6 - 3 = 3$.
\end{proposition}

The Ramanujan property here is notable: the collinearity graph achieves optimal expansion
for a 6-regular graph.  The four distinct eigenvalues $\{6, 3, -1, -3\}$ are all integers,
a signature of the graph's deep algebraic structure.

The multiplicity pattern $(1, 21, 27, 14)$ sums to 63 and reflects the representation
theory of $G_2(2)$.  Specifically, $G_2(2)$ acts on the 63 points; the decomposition of
the permutation character into irreducible representations of $G_2(2)$ yields exactly
these multiplicities.  The eigenvalue 3 with multiplicity 21 corresponds to the 21-dimensional
representation related to $\PSL(2,7)$; the eigenvalue $-1$ with multiplicity 27 and $-3$
with multiplicity 14 correspond to other irreducible components.

The girth of the collinearity graph is 3 (any three points on a common line form a triangle).
The graph is non-bipartite, and its spectrum is not symmetric about 0, consistent with
the eigenvalue $-1 \neq -6$.

\subsection{Meinhold-72 Graph}

The Meinhold-72 graph ($n = 72$, $k = 7$) has Ramanujan threshold $2\sqrt{6} \approx 4.899$.

\begin{table}[h]
\centering
\caption{Adjacency spectrum of the Meinhold-72 graph ($n=72$, $k=7$, girth 4).}
\label{tab:meinhold}
\begin{tabular}{@{}rrl@{}}
\toprule
Eigenvalue & Multiplicity & Exact form \\
\midrule
$7$           & $1$  & $k$ (trivial) \\
$\sqrt{13}$   & $14$ & $\approx 3.606$ \\
$1$           & $21$ & \\
$-1$          & $21$ & \\
$-\sqrt{13}$  & $14$ & \\
$-7$          & $1$  & \\
\bottomrule
\end{tabular}
\end{table}

\begin{proposition}
The Meinhold-72 graph has spectrum
\[
  \bigl\{7^1,\; (\sqrt{13})^{14},\; 1^{21},\; (-1)^{21},\;
  (-\sqrt{13})^{14},\; (-7)^1\bigr\}.
\]
Its spectrum is symmetric about 0, so it is bipartite.
It is \emph{not} Ramanujan: $\lambda_2 = \sqrt{13} \approx 3.606$ is within
$2\sqrt{6} \approx 4.899$, but $|\lambda_{\min}| = 7 = k$ exceeds it
(bipartite graphs always have $\lambda_{\min} = -k$).
\end{proposition}

The eigenvalue $-k = -7$ with multiplicity 1 is the hallmark of a bipartite regular graph:
$-\lambda$ is an eigenvalue whenever $\lambda$ is, and the trivial eigenvalue $k$ pairs with
$-k$.  The symmetric spectrum $(7, \sqrt{13}, 1, -1, -\sqrt{13}, -7)$ confirms bipartiteness.
Note that $13 = 7 + 6 = k + (k-1)$, an arithmetic relationship that sometimes occurs in
bipartite Cayley graphs.

The spectral gap $7 - \sqrt{13} \approx 3.394$ is the largest among the three graphs,
indicating strong connectivity despite the bipartite structure.

\section{Comparison and Discussion}

\subsection{Spectral Summary}

\begin{table}[h]
\centering
\caption{Spectral comparison of the three $G_2(2)$ graphs.}
\label{tab:comparison}
\begin{tabular}{@{}lrrp{3cm}p{1.2cm}ccrrrr@{}}
\toprule
Graph & $n$ & $k$ & $\lambda_2$ & $\lambda_{\min}$ & Ram.\ & Bip.\ & $g$ & $d$ & \#ev \\
\midrule
Tutte 12-cage        & 126 & 3 & $\sqrt{6}\approx 2.449$ & $-3$ & No  & Yes & 12 & 6 & 7 \\
GH(2,2) collinearity & 63  & 6 & $3$                     & $-3$ & Yes & No  &  3 & 3 & 4 \\
Meinhold-72          & 72  & 7 & $\sqrt{13}\approx 3.606$& $-7$ & No  & Yes &  4 & 5 & 6 \\
\bottomrule
\end{tabular}
\end{table}

\subsection{Same Group, Radically Different Spectra}

All three graphs have $\Aut(\Gamma) \cong G_2(2)$, yet their spectral properties differ
dramatically.  This illustrates a fundamental principle: the automorphism group constrains
but does not determine the spectrum.  The spectrum depends on the combinatorial structure of
the graph, not just its symmetry group.

The collinearity graph is the ``cleanest'' spectrally: four distinct integer eigenvalues,
Ramanujan, with a spectral gap equal to half the degree.  This reflects the fact that the
collinearity graph is a \emph{strongly regular} analogue (more precisely, distance-regular
with diameter 3), where the eigenvalues are determined by the intersection array and must
satisfy specific integrality conditions imposed by the Krein conditions.

The Tutte cage has the most eigenvalues (7, matching its diameter as required by
distance-regularity), and its irrational eigenvalues $\pm\sqrt{6}$ and $\pm\sqrt{2}$
arise from the characteristic polynomial of the distance-regular scheme, which factors
over $\mathbb{Q}(\sqrt{2}, \sqrt{3})$.

The Meinhold-72 graph sits between the other two: bipartite like the Tutte cage but
with fewer distinct eigenvalues (6), and with irrational $\pm\sqrt{13}$.  Its large
degree ($k = 7$) gives it the biggest spectral gap of the three.

\subsection{Ramanujan Analysis}

Only the collinearity graph is Ramanujan.  This is perhaps counterintuitive since the
Tutte cage is the most constrained (cubic, maximum girth) but fails to be Ramanujan.
The Alon--Boppana theorem guarantees that for any infinite family of $k$-regular graphs,
$\liminf \lambda_2 \geq 2\sqrt{k-1}$; the failure of the Tutte cage to be Ramanujan
shows that high girth does not automatically imply good expansion.

For the Meinhold-72 graph, even though $\lambda_2 = \sqrt{13} \approx 3.606 < 2\sqrt{6}
\approx 4.899$, the graph fails to be Ramanujan because $|\lambda_{\min}| = 7 = k$,
which is the maximum possible value.  This is again a consequence of bipartiteness:
bipartite $k$-regular graphs always have $\lambda_{\min} = -k$, so they can never be
Ramanujan in the strict sense.

\subsection{Distance-Regularity}

Two of the three graphs are distance-regular:
\begin{itemize}
  \item The \textbf{Tutte 12-cage} is distance-transitive (and hence distance-regular)
        with intersection array $\{3,2,2,2,2,2;1,1,1,1,1,3\}$ and diameter 6.
  \item The \textbf{$\GH$ collinearity graph} is distance-regular with diameter 3
        and intersection array $\{6,4,4;1,1,3\}$.
\end{itemize}
The \textbf{Meinhold-72 graph} is not distance-regular (diameter 5, but the number of
distinct eigenvalues is 6, and the intersection array does not exist globally).

A distance-regular graph of diameter $d$ has exactly $d+1$ distinct eigenvalues.
The Tutte cage: $d = 6$, $d+1 = 7$ distinct eigenvalues. \checkmark
The collinearity graph: $d = 3$, $d+1 = 4$ distinct eigenvalues. \checkmark

\subsection{Bipartiteness and Spectral Symmetry}

A $k$-regular graph is bipartite if and only if $-k$ is an eigenvalue.  Equivalently,
the spectrum is symmetric about 0.  The Tutte cage and Meinhold-72 are bipartite; the
collinearity graph is not (it contains triangles).

The bipartiteness of the Tutte cage is clear: points and lines form the two parts.
The bipartiteness of Meinhold-72 is more subtle and reflects the structure of the
Cayworth-to-Meinhold quotient construction.

\section{Algebraic Connectivity}

The \emph{algebraic connectivity} (Fiedler eigenvalue) of a connected graph is the
second-smallest eigenvalue of the Laplacian $L = kI - A$, equal to $k - \lambda_{\max}^-$
where $\lambda_{\max}^-$ is the largest eigenvalue strictly less than $k$.

\begin{table}[h]
\centering
\caption{Algebraic connectivity $\mu_2$ of the three graphs.}
\label{tab:fiedler}
\begin{tabular}{@{}lrr@{}}
\toprule
Graph & $\mu_2 = k - \lambda_2^*$ & Approximate value \\
\midrule
Tutte 12-cage        & $3 - \sqrt{6}$  & $\approx 0.551$ \\
GH(2,2) collinearity & $6 - 3 = 3$     & $3.000$ \\
Meinhold-72          & $7 - \sqrt{13}$ & $\approx 3.394$ \\
\bottomrule
\end{tabular}
\end{table}

The collinearity graph and Meinhold-72 both have large algebraic connectivity (3 and
$\approx 3.394$ respectively), meaning they are highly connected and difficult to
disconnect by removing few edges.  The Tutte cage, being sparse (cubic) and of large
girth, has much smaller algebraic connectivity $\approx 0.551$.

\section{Conclusion}

The three graphs arising from $G_2(2)$ present a case study in how a fixed automorphism
group can give rise to graphs with fundamentally different spectral behaviour.  The GH(2,2)
collinearity graph stands out as a particularly elegant object: a Ramanujan graph with
only four integer eigenvalues, optimal expansion, and deep connections to the representation
theory of the exceptional group $G_2(2)$ and its subgroups.

The Tutte 12-cage is the unique graph achieving the theoretical minimum size for a cubic
graph of girth 12, and its spectrum reflects the distance-regular structure of the
generalized hexagon.  The Meinhold-72 graph, constructed via a quotient of the Clayworth
graph, is the most mysterious of the three: bipartite, with irrational eigenvalues $\pm\sqrt{13}$,
and not distance-regular despite having $G_2(2)$ as its full automorphism group.

A natural direction for further investigation is whether these spectral profiles can be
predicted from the representation theory of $G_2(2)$: the eigenvalues of a Cayley or
coset graph are determined by the characters of the irreducible representations of the group,
evaluated at the connection set.  Matching the observed multiplicities $(1, 21, 27, 14)$
of the collinearity graph to the character table of $G_2(2)$ would give a representation-
theoretic proof of the spectral data.

\end{document}
